\documentclass[UTF8, 12pt, fontset=fandol, oneside]{ctexart}
\usepackage{researchnotes}

\title{第 10 章\quad 双线性型}
\author{}
\date{}

\begin{document}

\maketitle

\section*{第 10 章\quad 双线性型}

\subsection*{§10.1\quad 基本概念}

\subsubsection*{10.1.1\quad 对偶空间}

\noindent\textbf{1.\quad 对偶空间}

设 $V$ 是数域 $\mathbb{F}$ 上的线性空间，由 $V$ 到 $\mathbb{F}$ 上的线性映射（即线性函数）全体组成的线性空间 $V^*$ 称为 $V$ 的共轭空间. 当 $V$ 是有限维空间时，$V^*$ 称为 $V$ 的对偶空间.

\noindent\textbf{2.\quad 对偶基}

设 $V$ 是数域 $\mathbb{F}$ 上的 $n$ 维线性空间，$e_1,e_2,\cdots,e_n$ 是 $V$ 的一组基，$V$ 上的线性函数 $f_i$ 定义为 $f_i(e_i)=1,\ f_i(e_j)=0\ (j\ne i)$，则 $f_1,f_2,\cdots,f_n$ 是对偶空间 $V^*$ 的一组基，称为 $e_1,e_2,\cdots,e_n$ 的对偶基. 特别地，$\dim V^*=\dim V$.

\noindent\textbf{3.\quad 记号 $\langle\ ,\ \rangle$}

定义 $\langle f,x\rangle=f(x)$，其中 $f\in V^*,\ x\in V$，则 $\langle f,-\rangle=f$ 是 $V$ 上的线性函数，$\langle -,x\rangle$ 是 $V^*$ 上的线性函数. 定义线性映射 $\eta:V\to (V^*)^*=V^{**}$，$\eta(x)=\langle -,x\rangle$.

\noindent\textbf{4.\quad 定理}

当 $V$ 是有限维空间时，线性映射 $\eta:V\to V^{**}$ 是线性同构. 如果把 $V$ 与 $V^{**}$ 在这个同构下等同起来，则 $V$ 可以看成是 $V^*$ 的对偶空间，从而 $V$ 与 $V^*$ 互为对偶.

\noindent\textbf{5.\quad 定理}

设 $V,U$ 是数域 $\mathbb{F}$ 上的线性空间，$\varphi$ 是 $V$ 到 $U$ 的线性映射，则存在唯一的 $U^*$ 到 $V^*$ 的线性映射 $\varphi^*$，使得对任意的 $x\in V,\ f\in U^*$ 满足等式：
\begin{equation}
\langle \varphi^*(f),x\rangle=\langle f,\varphi(x)\rangle .
\tag{10.1}
\end{equation}
线性映射 $\varphi^*$ 称为 $\varphi$ 的对偶映射. 对偶映射具有下列性质：

(1) $(k_1\varphi_1+k_2\varphi_2)^*=k_1\varphi_1^*+k_2\varphi_2^*$，其中 $\varphi_1,\varphi_2\in\mathcal{L}(V,U),\ k_1,k_2\in\mathbb{F}$;

(2) $(\psi\varphi)^*=\varphi^*\psi^*$，其中 $\varphi\in\mathcal{L}(V,U),\ \psi\in\mathcal{L}(U,W)$;

(3) 若 $\varphi:V\to U$ 是线性同构，则 $\varphi^*:U^*\to V^*$ 也是线性同构，此时 $(\varphi^*)^{-1}=(\varphi^{-1})^*$.

\noindent\textbf{6.\quad 定理}

设 $V,U$ 是有限维线性空间，$\varphi:V\to U$ 是线性映射，$\varphi^*$ 是 $\varphi$ 的对偶映射.

(1) 设 $\{e_1,\cdots,e_n\}$ 是 $V$ 的一组基，$\{f_1,\cdots,f_n\}$ 是其对偶基；$\{u_1,\cdots,u_m\}$ 是 $U$ 的一组基，$\{g_1,\cdots,g_m\}$ 是其对偶基；$\varphi$ 在基 $\{e_1,\cdots,e_n\}$ 和基 $\{u_1,\cdots,u_m\}$ 下的表示矩阵是 $A$，则 $\varphi^*$ 在基 $\{g_1,\cdots,g_m\}$ 和基 $\{f_1,\cdots,f_n\}$ 下的表示矩阵是 $A'$.

(2) $\varphi$ 是单映射的充要条件是 $\varphi^*$ 是满映射，$\varphi$ 是满映射的充要条件是 $\varphi^*$ 是单映射. 特别地，$\varphi$ 是线性同构的充要条件是 $\varphi^*$ 也是线性同构.

\subsubsection*{10.1.2\quad 双线性型}

\noindent\textbf{1.\quad 双线性型}

设 $U,V$ 是数域 $\mathbb{F}$ 上的线性空间，$U\times V$ 是它们的积集合，若存在 $U\times V$ 到 $\mathbb{F}$ 的映射 $g$ 适合下列条件：

(1) 对任意的 $x,y\in U,\ z\in V,\ \lambda\in\mathbb{F}$,
\[
g(x+y,z)=g(x,z)+g(y,z),\quad g(\lambda x,z)=\lambda g(x,z);
\]

(2) 对任意的 $x\in U,\ z,w\in V,\ \lambda\in\mathbb{F}$,
\[
g(x,z+w)=g(x,z)+g(x,w),\quad g(x,\lambda z)=\lambda g(x,z),
\]
则称 $g$ 是 $U$ 和 $V$ 上的双线性函数或双线性型.

当 $U,V$ 是有限维线性空间时，任一 $U\times V$ 上的双线性型均可用矩阵来表示. 记 $\alpha_1,\alpha_2,\cdots,\alpha_m$ 是 $U$ 的基，$\beta_1,\beta_2,\cdots,\beta_n$ 是 $V$ 的基，令 $a_{ij}=g(\alpha_i,\beta_j)$，则
\[
G=\left(\begin{array}{cccc}
a_{11} & a_{12} & \cdots & a_{1n}\\
a_{21} & a_{22} & \cdots & a_{2n}\\
\vdots & \vdots &        & \vdots\\
a_{m1} & a_{m2} & \cdots & a_{mn}
\end{array}\right)
\]
称为 $g$ 在给定基下的表示矩阵. 设 $\alpha=x_1\alpha_1+x_2\alpha_2+\cdots+x_m\alpha_m,\ \beta=y_1\beta_1+y_2\beta_2+\cdots+y_n\beta_n$，$x=(x_1,x_2,\cdots,x_m)'$，$y=(y_1,y_2,\cdots,y_n)'$ 分别为 $\alpha,\beta$ 的坐标向量，则
\[
g(\alpha,\beta)=x'Gy.
\]
表示矩阵 $G$ 的秩称为双线性型 $g$ 的秩，记为 $\operatorname{r}(g)$.

\noindent\textbf{2.\quad 定理}

设 $g$ 是有限维线性空间 $U,V$ 上的双线性型，则总存在 $U,V$ 的基，使得 $g$ 在这两组基下的表示矩阵为相抵标准型
\[
\left(\begin{array}{cc}
I_r & O\\
O & O
\end{array}\right),
\]
其中 $r=\operatorname{r}(g)$.

\noindent\textbf{3.\quad 根子空间}

设 $g$ 是线性空间 $U,V$ 上的双线性型，令
\[
L=\{x\in U\mid g(x,y)=0\text{ 对一切 }y\in V\text{ 成立}\},
\]
\[
R=\{y\in V\mid g(x,y)=0\text{ 对一切 }x\in U\text{ 成立}\},
\]
则 $L$ 称为 $g$ 的左根子空间，$R$ 称为 $g$ 的右根子空间.

若 $g$ 的左、右根子空间都等于零，则称 $g$ 是非退化的双线性型.

\noindent\textbf{4.\quad 定理}

设 $g$ 是线性空间 $U,V$ 上的双线性型，则 $g$ 非退化的充要条件是
\[
\dim U=\dim V=\operatorname{r}(g).
\]
等价地，$g$ 非退化的充要条件是它的表示矩阵为可逆矩阵.

\noindent\textbf{5.\quad 定理}

设 $g_1,g_2$ 是线性空间 $U,V$ 上的两个非退化双线性型，则存在 $U$ 上的可逆线性变换 $\varphi$ 及 $V$ 上的可逆线性变换 $\psi$，使得对一切 $x\in U,\ y\in V$，有
\[
g_2(\varphi(x),y)=g_1(x,y),\quad g_2(x,\psi(y))=g_1(x,y).
\]

\subsubsection*{10.1.3\quad 纯量积}

\noindent\textbf{1.\quad 纯量积}

设 $g$ 是线性空间 $U=V$ 上的双线性型，称 $g$ 是 $V$ 上的一个纯量积（或数量积）.

\noindent\textbf{2.\quad 对称型和交错型}

设 $g$ 是线性空间 $V$ 上的纯量积，若对任意的 $x,y\in V$，都有
\[
g(x,y)=g(y,x),
\]
则称 $g$ 是 $V$ 上的对称型；若对任意的 $x,y\in V$，都有
\[
g(x,y)=-g(y,x),
\]
则称 $g$ 是 $V$ 上的交错型（或反对称型）.

\noindent\textbf{3.\quad 正交}

设 $g$ 是线性空间 $V$ 上的纯量积，$x,y\in V$，若 $g(x,y)=0$，则称 $x$ 左正交（或左垂直）于 $y$，称 $y$ 右正交（或右垂直）于 $x$，记为 $x\perp y$.

\noindent\textbf{4.\quad 定理}

设 $g$ 是线性空间 $V$ 上的纯量积，若对任意的 $x,y\in V$ 都有 $x\perp y$ 当且仅当 $y\perp x$，则 $g$ 必是对称型或交错型.

\noindent\textbf{5.\quad 定理}

设 $g_1,g_2$ 是线性空间 $V$ 上的非退化纯量积，则存在 $V$ 上唯一的可逆线性变换 $\varphi$，使得对任意的 $x,y\in V$，都有
\[
g_2(\varphi(x),y)=g_1(x,y).
\]

\subsubsection*{10.1.4\quad 交错型与辛空间}

\noindent\textbf{1.\quad 辛空间}

设 $V$ 是数域 $\mathbb{F}$ 上的线性空间，若在 $V$ 上定义了一个非退化的交错型，则称 $V$ 为辛空间.

\noindent\textbf{2.\quad 定理}

设 $g$ 是 $V$ 上的交错型，则存在 $V$ 的一组基，使得 $g$ 在这组基下的表示矩阵为分块对角矩阵：
\[
\operatorname{diag}\{S,\cdots,S,0,\cdots,0\},
\]
其中
\[
S=\left(\begin{array}{cc}
0 & 1\\
-1 & 0
\end{array}\right),
\]
这组基称为 $V$ 的辛基.

\noindent\textbf{3.\quad 辛变换}

设 $V$ 是辛空间，$\varphi$ 是 $V$ 上的可逆线性变换，若 $g(\varphi(x),\varphi(y))=g(x,y)$ 对任意的 $x,y\in V$ 成立，则称 $\varphi$ 是 $V$ 上的辛变换.

\noindent\textbf{4.\quad 定理}

设 $V$ 是数域 $\mathbb{F}$ 上的辛空间，则

(1) $V$ 上的线性变换 $\varphi$ 是辛变换的充要条件是 $\varphi$ 将辛基变到辛基；

(2) 两个辛变换之积仍是辛变换；

(3) 恒等变换是辛变换；

(4) 辛变换的逆变换是辛变换.

\subsubsection*{10.1.5\quad 对称型和正交空间}

\noindent\textbf{1.\quad 正交空间}

设 $V$ 是数域 $\mathbb{F}$ 上的线性空间，若在 $V$ 上定义了一个非退化的对称型，则称 $V$ 为（正则）正交空间.

\noindent\textbf{2.\quad 定理}

设 $g$ 是 $V$ 上的对称型，则必存在 $V$ 的一组基，使得 $g$ 在这组基下的表示矩阵为对角矩阵：
\[
\operatorname{diag}\{b_1,\cdots,b_r,0,\cdots,0\},
\]
这组基称为 $V$ 的正交基.

\noindent\textbf{3.\quad 迷向向量}

设 $x$ 是正交空间 $V$ 中的非零向量，若 $g(x,x)=0$，则称 $x$ 是迷向向量. 含有迷向向量的子空间称为迷向子空间.

\noindent\textbf{4.\quad 正交变换}

设 $V$ 是正交空间，$\varphi$ 是 $V$ 上的可逆线性变换，若 $g(\varphi(x),\varphi(y))=g(x,y)$ 对任意的 $x,y\in V$ 成立，则称 $\varphi$ 是 $V$ 上的正交变换.

\noindent\textbf{5.\quad 定理}

设 $V$ 是数域 $\mathbb{F}$ 上的正交空间，则

(1) 两个正交变换之积是正交变换；

(2) 恒等变换是正交变换；

(3) 正交变换的逆变换是正交变换.

\subsection*{§10.2\quad 线性函数与对偶空间}

线性空间的对偶空间是一个重要的概念，它在后续的专业课程以及物理学等领域中都有着广泛的应用. 通常的高等代数课程只讲授数域上的有限维线性空间理论，对无限维线性空间的情形涉及不多，比如一般并不给出无限维线性空间中基的定义及其存在性证明（这需要集合论中的选择公理或 Zorn 引理）. 因此除非特意指明，本章的大部分例题一般都在有限维线性空间的范畴内进行讨论. 例如，教材 [1] 给出了 §§10.1.1 定理 6 (2) 在有限维线性空间情形的证明，但只要建立了无限维线性空间中基的概念及其存在性，同样可证明 (2) 对无限维线性空间也成立. 然而，只有当 $V$ 是有限维线性空间时，才能由对偶基的存在性推出 $\dim V^*=\dim V$ 成立；当 $V$ 是无限维线性空间时，上述等式将不再成立，并且 §§10.1.1 定理 4 中的 $\eta:V\to V^{**}$ 也不再是线性同构. 由于这些结论的证明涉及到集合论和抽象代数的一些理论，故这里不准备展开阐述，有兴趣的读者可参考 [5].

\noindent\textbf{例 10.1}\quad 设 $V$ 是数域 $\mathbb{F}$ 上的线性空间（不必假设维数有限），$f,g$ 是 $V$ 上的非零线性函数，求证：$f$ 和 $g$ 线性相关的充要条件是 $\operatorname{Ker}f=\operatorname{Ker}g$.

\noindent\textbf{证明}\quad 若 $f=kg$，则显然 $\operatorname{Ker}f=\operatorname{Ker}g$. 下证充分性. 由 $f\ne0$ 可知，存在 $\alpha\in V$，使得 $f(\alpha)\ne0$，故可设 $g(\alpha)=kf(\alpha)$. 对任意的 $v\in V$，若设 $f(v)=cf(\alpha)$，则 $f(v-c\alpha)=0$，即 $v-c\alpha\in\operatorname{Ker}f=\operatorname{Ker}g$，从而 $g(v-c\alpha)=0$，故 $g(v)=cg(\alpha)$. 因此，对任意的 $v\in V$ 有
\[
g(v)=cg(\alpha)=ckf(\alpha)=kcf(\alpha)=kf(v),
\]
于是 $g=kf$，即 $f$ 和 $g$ 线性相关. $\square$

\noindent\textbf{例 10.2}\quad 设 $V$ 是数域 $\mathbb{F}$ 上的 $n$ 维线性空间，$f,g$ 是 $V$ 上的非零线性函数. 求证：若 $f,g$ 线性无关，则对任意的 $v\in V$，存在分解 $v=u+w$，使得 $f(v)=f(w)$，$g(v)=g(u)$.

\noindent\textbf{证明}\quad 设 $e_1,e_2,\cdots,e_n$ 是 $V$ 的一组基，则 $\alpha=\sum_{i=1}^n c_ie_i\in\operatorname{Ker}f$ 当且仅当
\[
0=f(\alpha)=f\left(\sum_{i=1}^n c_ie_i\right)=\sum_{i=1}^n c_if(e_i).
\]
换言之，$\alpha\in\operatorname{Ker}f$ 当且仅当 $\alpha$ 的坐标向量 $(c_1,c_2,\cdots,c_n)'$ 是线性方程 $f(e_1)x_1+f(e_2)x_2+\cdots+f(e_n)x_n=0$ 的解. 由于 $f,g$ 都是非零线性函数，故由线性映射的维数公式可知 $\dim\operatorname{Ker}f=n-1,\ \dim\operatorname{Ker}g=n-1$. 根据一开始的说明可知，$\operatorname{Ker}f\cap\operatorname{Ker}g$ 是下列联立线性方程组的解空间：
\[
\left\{
\begin{array}{l}
f(e_1)x_1+f(e_2)x_2+\cdots+f(e_n)x_n=0,\\
g(e_1)x_1+g(e_2)x_2+\cdots+g(e_n)x_n=0.
\end{array}
\right.
\]
若上述方程组的系数矩阵的秩等于 1，则存在 $k\in\mathbb{F}$，使得 $g(e_i)=kf(e_i)\ (1\le i\le n)$，于是 $g=kf$，这与 $f,g$ 线性无关矛盾. 因此上述方程组的系数矩阵的秩等于 2，从而 $\dim(\operatorname{Ker}f\cap\operatorname{Ker}g)=n-2$. 再由交和空间的维数公式可知
\[
\begin{aligned}
\dim(\operatorname{Ker}f+\operatorname{Ker}g)
&=\dim\operatorname{Ker}f+\dim\operatorname{Ker}g-\dim(\operatorname{Ker}f\cap\operatorname{Ker}g)\\
&=(n-1)+(n-1)-(n-2)=n=\dim V,
\end{aligned}
\]
于是 $V=\operatorname{Ker}f+\operatorname{Ker}g$. 因此对任意的 $v\in V$，存在分解 $v=u+w$，其中 $u\in\operatorname{Ker}f,\ w\in\operatorname{Ker}g$，使得 $f(v)=f(w),\ g(v)=g(u)$. $\square$

\noindent\textbf{注}\quad 由例 10.2 的证明方法不难得到例 10.1 在有限维线性空间情形的另一证明，请读者自行补充完整.

\noindent\textbf{例 10.3}\quad 设 $V$ 是数域 $\mathbb{F}$ 上的 $n$ 维线性空间，$U$ 是 $V$ 的非平凡子空间，求证：必存在 $V$ 上的线性函数 $f_i(1\le i\le r)$，使得 $U=\bigcap_{i=1}^r\operatorname{Ker}f_i$.

\noindent\textbf{证明}\quad 设 $e_{r+1},\cdots,e_n$ 是 $U$ 的一组基，将它扩张为 $V$ 的一组基 $e_1,\cdots,e_r,e_{r+1},\cdots,e_n$. 设 $f_1,f_2,\cdots,f_n$ 为上述基的对偶基，即满足 $f_i(e_j)=\delta_{ij}$，则不难验证 $\operatorname{Ker}f_i=L(e_1,\cdots,e_{i-1},e_{i+1},\cdots,e_n)$，于是 $U=L(e_{r+1},\cdots,e_n)=\bigcap_{i=1}^r\operatorname{Ker}f_i$. $\square$

\noindent\textbf{例 10.4}\quad 设 $U,V$ 是数域 $\mathbb{F}$ 上的线性空间（不必假设维数有限），$U^*,V^*$ 分别是它们的共轭空间. 求证：
\[
U^*\oplus V^*\cong (U\oplus V)^*.
\]

\noindent\textbf{证明}\quad 设 $f_1\in U^*,\ f_2\in V^*$，定义 $f$ 为 $U\oplus V$ 上的线性函数：
\[
f(x+y)=f_1(x)+f_2(y),\quad x\in U,\ y\in V.
\]
令 $\varphi(f_1+f_2)=f$，则不难验证 $\varphi$ 是 $U^*\oplus V^*\to (U\oplus V)^*$ 的线性映射. 另一方面，假设 $f$ 是 $U\oplus V$ 上的线性函数，令 $f_1,f_2$ 分别是 $f$ 在 $U,V$ 上的限制，定义 $\psi$ 是 $(U\oplus V)^*\to U^*\oplus V^*$ 的线性映射：$\psi(f)=f_1+f_2$. 容易验证 $\psi\varphi$ 和 $\varphi\psi$ 分别是 $U^*\oplus V^*$ 和 $(U\oplus V)^*$ 上的恒等映射，因此 $\varphi$ 是线性同构. $\square$

\noindent\textbf{例 10.5}\quad 设 $V_1$ 是线性空间 $V$（不必假设维数有限）的子空间，记
\[
V_1^\perp=\{f\in V^*\mid \langle f,V_1\rangle=0\}.
\]
求证：$V_1^\perp$ 是 $V^*$ 的子空间，且若 $V_2$ 是 $V$ 的另外一个子空间，则
\[
V_1^\perp\cap V_2^\perp=(V_1+V_2)^\perp.
\]

\noindent\textbf{证明}\quad 容易验证 $V_1^\perp$ 是子空间. 若 $U,W$ 是 $V$ 的子空间且 $U\subseteq W$，显然有 $W^\perp\subseteq U^\perp$. 因此 $(V_1+V_2)^\perp\subseteq V_1^\perp,\ (V_1+V_2)^\perp\subseteq V_2^\perp$，从而 $(V_1+V_2)^\perp\subseteq V_1^\perp\cap V_2^\perp$. 反之，若 $f\in V_1^\perp\cap V_2^\perp$，则对任意的 $v_1\in V_1,\ v_2\in V_2$，$\langle f,v_1+v_2\rangle=f(v_1)+f(v_2)=0$，因此 $f\in(V_1+V_2)^\perp$，即有 $V_1^\perp\cap V_2^\perp\subseteq (V_1+V_2)^\perp$. 这就证明了后一个结论. $\square$

\noindent\textbf{例 10.6}\quad 设 $V$ 是数域 $\mathbb{F}$ 上的 $n$ 维线性空间，$V_1$ 是 $V$ 的子空间，求证：
\[
\dim V=\dim V_1+\dim V_1^\perp.
\]

\noindent\textbf{证明}\quad 取 $V_1$ 的一组基 $e_1,\cdots,e_r$，并扩张为 $V$ 的一组基 $e_1,e_2,\cdots,e_n$，再取其对偶基 $f_1,f_2,\cdots,f_n$. 由对偶基的定义可知 $f_j(e_i)=0\ (1\le i\le r,\ r+1\le j\le n)$，从而 $f_j(V_1)=0$，即 $f_j\in V_1^\perp\ (r+1\le j\le n)$. 另一方面，任取 $f\in V_1^\perp$，设 $f=a_1f_1+a_2f_2+\cdots+a_nf_n$，依次作用上 $e_1,\cdots,e_r$ 可得 $a_1=\cdots=a_r=0$，故 $f$ 是 $f_{r+1},\cdots,f_n$ 的线性组合. 因此 $f_{r+1},\cdots,f_n$ 是 $V_1^\perp$ 的一组基，特别地，$\dim V_1^\perp=n-r$，故结论成立. $\square$

\noindent\textbf{例 10.7}\quad 设 $V_1,V_2$ 是 $n$ 维线性空间 $V$ 的子空间，将 $V$ 看成是 $V^*$ 的对偶空间. 求证：
\[
(V_1^\perp)^\perp=V_1,\quad (V_1\cap V_2)^\perp=V_1^\perp+V_2^\perp.
\]

\noindent\textbf{证明}\quad 显然 $V_1\subseteq (V_1^\perp)^\perp$. 由例 10.6 可知 $\dim V_1^\perp=n-\dim V_1$，故 $\dim (V_1^\perp)^\perp=n-\dim V_1^\perp=\dim V_1$，于是 $(V_1^\perp)^\perp=V_1$. 由例 10.5 和第一个结论可知，
\[
(V_1^\perp+V_2^\perp)^\perp=(V_1^\perp)^\perp\cap (V_2^\perp)^\perp=V_1\cap V_2,
\]
再次由第一个结论可得 $(V_1\cap V_2)^\perp=((V_1^\perp+V_2^\perp)^\perp)^\perp=V_1^\perp+V_2^\perp$. $\square$

\noindent\textbf{例 10.8}\quad 设 $\varphi$ 是 $n$ 维线性空间 $V$ 上的线性变换，$\varphi^*$ 是 $\varphi$ 的对偶变换，求证：
\[
\operatorname{Im}\varphi^*=(\operatorname{Ker}\varphi)^\perp.
\]

\noindent\textbf{证法 1}\quad 假设 $f\in\operatorname{Im}\varphi^*$，则存在 $g\in V^*$，使得 $f=\varphi^*(g)$. 对 $\operatorname{Ker}\varphi$ 中任一向量 $x$，有
\[
\langle f,x\rangle=\langle\varphi^*(g),x\rangle=\langle g,\varphi(x)\rangle=0.
\]
因此 $f\in(\operatorname{Ker}\varphi)^\perp$，从而 $\operatorname{Im}\varphi^*\subseteq(\operatorname{Ker}\varphi)^\perp$.

另一方面，设 $\dim\operatorname{Ker}\varphi=k$，则由例 10.6 可得 $\dim(\operatorname{Ker}\varphi)^\perp=n-k$. 设 $\varphi$ 在 $V$ 的一组基 $\{e_1,\cdots,e_n\}$ 下的表示矩阵为 $A$，则 $\varphi^*$ 在 $V^*$ 的对偶基 $\{f_1,\cdots,f_n\}$ 下的表示矩阵为 $A'$. 于是 $\dim\operatorname{Im}\varphi^*=\operatorname{r}(A')=\operatorname{r}(A)=\dim\operatorname{Im}\varphi=n-k$，从而可得 $\operatorname{Im}\varphi^*=(\operatorname{Ker}\varphi)^\perp$.

\noindent\textbf{证法 2}\quad 由例 10.7 可知，我们只要证明 $\operatorname{Ker}\varphi=(\operatorname{Im}\varphi^*)^\perp$ 即可. 若 $x\in\operatorname{Ker}\varphi$，则对任意的 $\varphi^*(f)\in\operatorname{Im}\varphi^*$，有 $\langle\varphi^*(f),x\rangle=\langle f,\varphi(x)\rangle=0$，因此 $x\in(\operatorname{Im}\varphi^*)^\perp$，即 $\operatorname{Ker}\varphi\subseteq(\operatorname{Im}\varphi^*)^\perp$. 另一方面，任取 $x\in(\operatorname{Im}\varphi^*)^\perp$，则对任意的 $\varphi^*(f)\in\operatorname{Im}\varphi^*$，有 $0=\langle\varphi^*(f),x\rangle=\langle f,\varphi(x)\rangle$. 由 $f$ 的任意性可知 $\varphi(x)=0$，即 $x\in\operatorname{Ker}\varphi$，从而 $(\operatorname{Im}\varphi^*)^\perp\subseteq\operatorname{Ker}\varphi$，于是结论得证. $\square$

\noindent\textbf{注}\quad 例 10.8 证法 2 的好处是，证明 $\operatorname{Ker}\varphi=(\operatorname{Im}\varphi^*)^\perp$ 的过程不涉及维数的有限性，从而这一结论在无限维线性空间的情形依然成立（此时需要无限维线性空间基的存在性）. 然而 $\operatorname{Im}\varphi^*=(\operatorname{Ker}\varphi)^\perp$ 这一结论一般不能推广到无限维线性空间的情形，但在一些特殊情况下可以推广，我们来看下面的例题.

\noindent\textbf{例 10.9}\quad 设 $\varphi$ 是线性空间 $V$（不要求是有限维）上的幂等线性变换（即 $\varphi^2=\varphi$），$\varphi^*$ 是 $\varphi$ 的对偶变换，求证：
\[
\operatorname{Im}\varphi^*=(\operatorname{Ker}\varphi)^\perp.
\]

\noindent\textbf{证明}\quad 与例 10.8 完全一样的证明可得 $\operatorname{Im}\varphi^*\subseteq(\operatorname{Ker}\varphi)^\perp$. 另一方面，任取 $f\in(\operatorname{Ker}\varphi)^\perp$，如果能证明 $f=\varphi^*(f)$，就能得到 $f\in\operatorname{Im}\varphi^*$，从而 $\operatorname{Im}\varphi^*=(\operatorname{Ker}\varphi)^\perp$ 成立. 事实上，对任意的 $v\in V$，由 $\varphi^2=\varphi$ 可知 $v-\varphi(v)\in\operatorname{Ker}\varphi$，于是 $f(v-\varphi(v))=0$，从而 $f(v)=f(\varphi(v))=\varphi^*(f)(v)$ 对任意的 $v\in V$ 成立，因此 $f=\varphi^*(f)$. $\square$

\noindent\textbf{例 10.10}\quad 设 $\varphi$ 是 $n$ 维线性空间 $V$ 上的线性变换，$V_1$ 是 $V$ 的子空间，求证：$V_1$ 是 $\varphi$ 的不变子空间的充要条件是 $V_1^\perp$ 是 $\varphi^*$ 的不变子空间.

\noindent\textbf{证明}\quad 若 $V_1$ 是 $\varphi$ 的不变子空间，则对任意的 $v\in V_1$，有 $\varphi(v)\in V_1$，从而对任意的 $f\in V_1^\perp$，有 $\langle\varphi^*(f),v\rangle=\langle f,\varphi(v)\rangle=0$，即 $\varphi^*(f)\in V_1^\perp$，于是 $V_1^\perp$ 是 $\varphi^*$ 的不变子空间. 反之，若 $V_1^\perp$ 是 $\varphi^*$ 的不变子空间，则对任意的 $f\in V_1^\perp$，有 $\varphi^*(f)\in V_1^\perp$，从而对任意的 $v\in V_1$，有 $\langle f,\varphi(v)\rangle=\langle\varphi^*(f),v\rangle=0$，即 $\varphi(v)\in(V_1^\perp)^\perp=V_1$，于是 $V_1$ 是 $\varphi$ 的不变子空间. $\square$

\noindent\textbf{注}\quad 设 $V$ 是线性空间（不要求是有限维），$V_1$ 是 $V$ 的子空间，若承认无限维线性空间基的存在性，则可证明对任一 $v\notin V_1$，存在 $f\in V_1^\perp$，使得 $f(v)\ne0$. 如果有了这一结论，则例 10.10 的结论对无限维线性空间也成立.

\noindent\textbf{例 10.11}\quad 设 $V,U$ 是 $\mathbb{F}$ 上的有限维线性空间，$\varphi$ 是 $V\to U$ 的线性映射. 求证：若将 $V$ 与 $V^*$，$U$ 与 $U^*$ 看成是互为对偶的空间，则 $(\varphi^*)^*=\varphi$.

\noindent\textbf{证明}\quad 对任意的 $x\in V=V^{**},\ f\in U^*$，我们有
\[
\langle f,\varphi(x)\rangle=\langle\varphi^*(f),x\rangle=\langle f,\varphi^{**}(x)\rangle,
\]
因此 $\varphi(x)=\varphi^{**}(x)$，即 $\varphi=\varphi^{**}$. $\square$

\noindent\textbf{例 10.12}\quad 设 $V$ 是 $n$ 维欧氏空间，则对任一固定的 $u\in V$，$(u,-)$ 是 $V$ 上的线性函数，作映射 $\eta:V\to V^*,\ \eta(u)=(u,-)$. 证明：

(1) $\eta$ 是线性同构，特别地，若将 $u$ 与 $(u,-)$ 等同起来，则 $\langle u,v\rangle=(u,v)$，即可将 $V$ 看成是自身的对偶空间；

(2) $V$ 的任一组标准正交基 $e_1,e_2,\cdots,e_n$ 的对偶基是其自身；

(3) $V$ 上任一线性变换 $\varphi$ 的对偶变换就是 $\varphi$ 的伴随.

\noindent\textbf{证明}\quad (1) 容易验证 $(u,-)$ 是线性函数以及 $\eta$ 是线性映射. 假设 $\eta(u)=0$，则对任意的 $v\in V$，$(u,v)=\eta(u)(v)=0$，由内积的正定性可得 $u=0$，因此 $\eta$ 是单映射. 又 $\dim V^*=\dim V=n$，故由线性映射的维数公式可知 $\eta$ 是线性同构.

(2) 由 (1) 以及 $(e_i,e_j)=\delta_{ij}$ 即得结论.

(3) 记 $\varphi^\sharp$ 是 $\varphi$ 的对偶变换，$\varphi^*$ 是 $\varphi$ 的伴随，则由 (1) 可知
\[
(\varphi^\sharp(u),v)=\langle u,\varphi(v)\rangle=(u,\varphi(v))=(\varphi^*(u),v)=\langle\varphi^*(u),v\rangle,
\]
再由 $u,v$ 的任意性即得 $\varphi^\sharp=\varphi^*$. $\square$

\subsection*{§10.3\quad 双线性型与纯量积}

\noindent\textbf{例 10.13}\quad 设 $g$ 是 $U,V$ 上的非退化双线性型，若 $\{u_i\},\{v_i\}(1\le i\le n)$ 分别是 $U,V$ 的基，使得 $g(u_i,v_j)=\delta_{ij}$，则称 $\{u_i\},\{v_i\}$ 是关于 $g$ 的对偶基. 设 $\varphi$ 是 $V$ 上
的线性变换，$\varphi^*$ 是 $\varphi$ 关于 $g$ 的对偶变换. 若 $\varphi$ 在基 $\{v_i\}$ 下的表示矩阵为 $A$，求证：$\varphi^*$ 在基 $\{u_i\}$ 下的表示矩阵是 $A'$.

\noindent\textbf{证明}\quad 设 $\varphi^*$ 在基 $\{u_i\}$ 下的表示矩阵为 $B=(b_{ij})$，则有
\[
\varphi^*(u_i)=b_{1i}u_1+b_{2i}u_2+\cdots+b_{ni}u_n,\quad 1\le i\le n.
\]
设 $A=(a_{ij})$，则有
\[
\varphi(v_j)=a_{1j}v_1+a_{2j}v_2+\cdots+a_{nj}v_n,\quad 1\le j\le n.
\]
注意到
\[
b_{ji}=g(\varphi^*(u_i),v_j)=g(u_i,\varphi(v_j))=a_{ij},\quad 1\le i,j\le n,
\]
此即 $B=A'$. $\square$

\noindent\textbf{例 10.14}\quad 设 $g$ 是 $U,V$ 上的非零双线性型，证明：必存在 $U,V$ 的子空间 $U_0,V_0$，使得 $g$ 在 $U_0,V_0$ 上的限制是非退化的双线性型，且
\[
\dim U_0=\dim V_0=\dim U-\dim L,
\]
其中 $L$ 是 $g$ 的左根子空间.

\noindent\textbf{证明}\quad 设 $g$ 在 $U$ 的基 $\{u_1,u_2,\cdots,u_m\}$ 和 $V$ 的基 $\{v_1,v_2,\cdots,v_n\}$ 下的表示矩阵为相抵标准型，即
\[
A=\left(\begin{array}{cc}
I_r & O\\
O & O
\end{array}\right).
\]
令 $U_0$ 是由基向量 $u_1,\cdots,u_r$ 生成的子空间，$V_0$ 是由基向量 $v_1,\cdots,v_r$ 生成的子空间. 显然，将 $g$ 限制在 $U_0,V_0$ 上是非退化的双线性型，且 $\dim U_0=\dim V_0=r$. 又 $g$ 的左根子空间就是由 $u_{r+1},\cdots,u_m$ 生成的子空间，因此 $\dim L=m-r$，即有 $\dim U_0=\dim U-\dim L$. $\square$

\noindent\textbf{例 10.15}\quad 设 $g$ 是 $U,V$ 上的非退化双线性型，$\varphi,\psi$ 是 $V$ 上的线性变换，求证：

(1) $(k\varphi+l\psi)^*=k\varphi^*+l\psi^*$，其中 $k,l$ 是常数；

(2) $(\psi\varphi)^*=\varphi^*\psi^*$;

(3) 若 $\varphi$ 是 $V$ 的自同构，则 $\varphi^*$ 是 $U$ 的自同构，此时 $(\varphi^*)^{-1}=(\varphi^{-1})^*$;

(4) $(\varphi^*)^*=\varphi$.

\noindent\textbf{证明}\quad (1) 对任意的 $u\in U,\ v\in V$，由对偶变换的定义可得
\[
\begin{aligned}
g(u,(k\varphi+l\psi)(v))
&=g(u,k\varphi(v))+g(u,l\psi(v))\\
&=g(k\varphi^*(u),v)+g(l\psi^*(u),v)\\
&=g((k\varphi+l\psi)^*(u),v),
\end{aligned}
\]
再由对偶变换的唯一性即得 $(k\varphi+l\psi)^*=k\varphi^*+l\psi^*$.

(2) 对任意的 $u\in U,\ v\in V$，由对偶变换的定义可得
\[
g(u,\psi\varphi(v))=g(\psi^*(u),\varphi(v))=g(\varphi^*\psi^*(u),v),
\]
再由对偶变换的唯一性即得 $(\psi\varphi)^*=\varphi^*\psi^*$.

(3) 若 $\varphi$ 是 $V$ 的自同构，则 $\varphi^{-1}\varphi=\varphi\varphi^{-1}=I_V$. 两边同取对偶，由 (2) 可得
\[
\varphi^*(\varphi^{-1})^*=(\varphi^{-1})^*\varphi^*=I_V^*=I_U,
\]
故 $\varphi^*$ 是 $U$ 的自同构，并且 $(\varphi^*)^{-1}=(\varphi^{-1})^*$.

(4) 对任意的 $u\in U,\ v\in V$，由对偶变换的定义可得
\[
g(u,\varphi(v))=g(\varphi^*(u),v)=g(u,(\varphi^*)^*(v)),
\]
再由对偶变换的唯一性即得 $(\varphi^*)^*=\varphi$. 本题也可利用例 10.13 的结论来证明. $\square$

\noindent\textbf{例 10.16}\quad 设 $g,h$ 是 $n$ 维线性空间 $V$ 上秩相同的纯量积，求证：必存在 $V$ 上的可逆线性变换 $\varphi,\psi$，使得 $h(x,y)=g(\varphi(x),\psi(y))$ 对一切 $x,y\in V$ 成立.

\noindent\textbf{证明}\quad 我们用矩阵方法来证明结论. 设 $g,h$ 在 $V$ 的某一组基下的表示矩阵分别为 $A,B$，向量 $x,y$ 的坐标向量（用列向量表示）分别为 $\alpha,\beta$，则
\[
g(x,y)=\alpha'A\beta,\quad h(x,y)=\alpha'B\beta.
\]
又假设线性变换 $\varphi$ 和 $\psi$ 在同一组基下的表示矩阵分别为 $C,D$（待定），则
\[
g(\varphi(x),\psi(y))=(C\alpha)'A(D\beta)=\alpha'C'AD\beta.
\]
因为 $A$ 和 $B$ 秩相同，故存在可逆矩阵 $C,D$，使得 $C'AD=B$，于是结论得证. $\square$

\noindent\textbf{例 10.17}\quad 设 $W=U\oplus V$，$g$ 是 $U$ 上的纯量积，$h$ 是 $V$ 上的纯量积. 现定义 $W$ 上的纯量积 $q$ 如下：
\[
q(x+y,u+v)=g(x,u)+h(y,v),
\]
其中 $x,u\in U,\ y,v\in V$，求证：

(1) 若 $g,h$ 非退化，则 $q$ 也非退化；

(2) 若 $g,h$ 是对称型（交错型），则 $q$ 也是对称型（交错型）；

(3) 若 $\{u_i\},\{v_i\}$ 分别是 $U$ 和 $V$ 的基，且 $g,h$ 在这两组基下的表示矩阵分别为 $A,B$，则 $q$ 在 $W$ 的基 $\{u_i\}\cup\{v_i\}$ 下的表示矩阵为分块对角矩阵 $\operatorname{diag}\{A,B\}$.

\noindent\textbf{证明}\quad 若矩阵 $A$ 和 $B$ 可逆，则 $\operatorname{diag}\{A,B\}$ 也可逆，因此 (1) 是 (3) 的推论. (2) 的验证很容易，现只需证明 (3). 因为
\[
q(u_i,v_j)=q(u_i+0,0+v_j)=g(u_i,0)+h(0,v_j)=0,
\]
以及 $q(v_j,u_i)=0$，所以 $q$ 的表示矩阵是分块对角矩阵. 又
\[
q(u_i,u_j)=g(u_i,u_j),\quad q(v_i,v_j)=h(v_i,v_j),
\]
因此结论成立. $\square$

\noindent\textbf{例 10.18}\quad 设 $V$ 是由 $n$ 阶实矩阵全体构成的欧氏空间（取 Frobenius 内积），则 Frobenius 内积 $(-,-)$ 是 $V$ 上的非退化对称型. 设 $A_1,\cdots,A_{n^2}$ 是 $V$ 的一组基，$B_1,\cdots,B_{n^2}$ 是其对偶基，即满足 $(A_i,B_j)=\delta_{ij}\ (1\le i,j\le n^2)$. 求证：
\[
\sum_{i=1}^{n^2}A_iB_i=I_n.
\]

\noindent\textbf{证明}\quad 设 $E_{11},\cdots,E_{nn}$ 是 $n$ 阶基础矩阵，为书写方便将它们依次标记为 $E_1,\cdots,E_{n^2}$. 显然，这是 $V$ 的一组标准正交基，从而它的对偶基也是其自身. 设
\[
(A_1,\cdots,A_{n^2})=(E_1,\cdots,E_{n^2})P,\quad
(B_1,\cdots,B_{n^2})=(E_1,\cdots,E_{n^2})Q,
\]
其中 $P=(p_{ij}),\ Q=(q_{ij})$ 是基之间的过渡矩阵，则 $A_i=\sum_{k=1}^{n^2}p_{ki}E_k,\ B_j=\sum_{l=1}^{n^2}q_{lj}E_l$. 于是对任意的 $1\le i,j\le n^2$ 有
\[
\delta_{ij}=(A_i,B_j)=\left(\sum_{k=1}^{n^2}p_{ki}E_k,\sum_{l=1}^{n^2}q_{lj}E_l\right)=\sum_{k=1}^{n^2}p_{ki}q_{kj},
\]
这即为 $P'Q=I_{n^2}$. 于是 $QP'=I_{n^2}$，从而 $PQ'=I_{n^2}$，此即 $\sum_{i=1}^{n^2}p_{ki}q_{li}=\delta_{kl}$. 因此
\[
\begin{aligned}
\sum_{i=1}^{n^2}A_iB_i
&=\sum_{i,k,l=1}^{n^2}p_{ki}q_{li}E_kE_l\\
&=\sum_{k,l=1}^{n^2}\left(\sum_{i=1}^{n^2}p_{ki}q_{li}\right)E_kE_l\\
&=\sum_{k,l=1}^{n^2}\delta_{kl}E_kE_l
=\sum_{k=1}^{n^2}E_k^2\\
&=\sum_{i,j=1}^n E_{ij}E_{ij}
=\sum_{i=1}^n E_{ii}=I_n.
\end{aligned}
\]
$\square$

\noindent\textbf{例 10.19}\quad 设 $g$ 是 $n$ 维线性空间 $V$ 上的对称型或交错型，$U$ 是 $V$ 的子空间，求证：$U\cap U^\perp=0$ 的充要条件是 $g$ 限制在 $U$ 上是一个非退化的纯量积，这时有直和分解 $V=U\oplus U^\perp$.

\noindent\textbf{证明}\quad 设 $U\cap U^\perp=0$，若 $g$ 限制在 $U$ 上退化，则存在 $U$ 中非零向量 $u$，使得 $g(u,U)=0$，从而 $u\in U\cap U^\perp$，推出矛盾. 反之，设 $g$ 限制在 $U$ 上非退化，任取 $u\in U\cap U^\perp$，则 $g(u,U)=0$，从而 $u=0$，这表明 $U\cap U^\perp=0$.

对于第二个结论，我们先证明若 $g$ 限制在 $U$ 上非退化，则 $\dim U+\dim U^\perp=n$. 对任意的 $v\in V$，$g(v,-)$ 限制在 $U$ 上是 $U$ 上的线性函数. 作线性映射 $\varphi:V\to U^*$，$\varphi(v)=g(v,-)$，则 $\operatorname{Ker}\varphi=U^\perp$. 因为 $g$ 限制在 $U$ 上非退化，故限制映射 $\varphi|_U:U\to U^*$ 是单射，又 $\dim U=\dim U^*$，从而 $\varphi|_U:U\to U^*$ 是同构. 因此对任意的 $f\in U^*$，存在 $u\in U$，使得 $f=g(u,-)$，于是 $\varphi$ 是满射. 最后，由线性映射的维数公式即得 $\dim U+\dim U^\perp=n$，又因为 $U\cap U^\perp=0$，所以 $V=U\oplus U^\perp$. $\square$

\subsection*{§10.4\quad 交错型与辛几何}

\noindent\textbf{例 10.20}\quad 设 $h$ 是三维线性空间 $V$ 上的非零交错型，求证：存在 $V$ 上的线性函数 $f,g$，使得对任意的 $x,y\in V$，有 $h(x,y)=f(x)g(y)-f(y)g(x)$.

\noindent\textbf{证明}\quad 设 $h$ 在 $V$ 的基 $v_1,v_2,v_3$ 下的表示矩阵为标准型
\[
\left(\begin{array}{ccc}
0 & 1 & 0\\
-1 & 0 & 0\\
0 & 0 & 0
\end{array}\right),
\]
令
\[
f(x)=h(x,v_2),\quad g(x)=h(v_1,x),
\]
不难验证 $f,g$ 即为要求之线性函数. $\square$

我们通过下面这道例题来看一看如何求出辛空间的一组辛基.

\noindent\textbf{例 10.21}\quad 设四维辛空间 $(V,g)$ 在一组基 $\{e_1,e_2,e_3,e_4\}$ 下的表示矩阵为
\[
A=\left(\begin{array}{rrrr}
0 & 2 & -1 & 3\\
-2 & 0 & 4 & -2\\
1 & -4 & 0 & 1\\
-3 & 2 & -1 & 0
\end{array}\right),
\]
求 $V$ 的一组辛基.

\noindent\textbf{解}\quad 由例 8.16 可知，存在可逆矩阵 $C$，使得 $C'AC=B=\operatorname{diag}\{S,S\}$，其中
\[
S=\left(\begin{array}{cc}
0 & 1\\
-1 & 0
\end{array}\right).
\]
基之间的过渡矩阵 $C$ 可用对称初等变换法来求（类似于对称矩阵合同于对角矩阵的求法）：对矩阵 $(A\vdots I)$ 施以初等行变换，再对 $A$ 施以对称的初等列变换，直到将 $A$ 变为 $B$，此时 $I$ 就变成 $C'$，转置后即得 $C$.

将 $\dfrac12$ 乘以 $(A\vdots I)$ 的第二行，再乘以第二列得到
\[
(A\vdots I)\longrightarrow
\left(\begin{array}{rrrr|rrrr}
0 & 1 & -1 & 3 & 1 & 0 & 0 & 0\\
-1 & 0 & 2 & -1 & 0 & \frac12 & 0 & 0\\
1 & -2 & 0 & 1 & 0 & 0 & 1 & 0\\
-3 & 1 & -1 & 0 & 0 & 0 & 0 & 1
\end{array}\right).
\]
将第二行分别乘以 $1,-3$ 后加到第三行及第四行上，再将第二列分别乘以 $1,-3$ 后加到第三列及第四列上得到
\[
\longrightarrow
\left(\begin{array}{rrrr|rrrr}
0 & 1 & 0 & 0 & 1 & 0 & 0 & 0\\
-1 & 0 & 2 & -1 & 0 & \frac12 & 0 & 0\\
0 & -2 & 0 & 6 & 0 & \frac12 & 1 & 0\\
0 & 1 & -6 & 0 & 0 & -\frac32 & 0 & 1
\end{array}\right).
\]
将第一行分别乘以 $2,-1$ 后加到第三行及第四行上，再将第一列分别乘以 $2,-1$ 后加到第三列及第四列上得到
\[
\longrightarrow
\left(\begin{array}{rrrr|rrrr}
0 & 1 & 0 & 0 & 1 & 0 & 0 & 0\\
-1 & 0 & 0 & 0 & 0 & \frac12 & 0 & 0\\
0 & 0 & 0 & 6 & 2 & \frac12 & 1 & 0\\
0 & 0 & -6 & 0 & -1 & -\frac32 & 0 & 1
\end{array}\right).
\]
将第四行乘以 $\dfrac16$，再将第四列乘以 $\dfrac16$ 后得到
\[
\longrightarrow
\left(\begin{array}{rrrr|rrrr}
0 & 1 & 0 & 0 & 1 & 0 & 0 & 0\\
-1 & 0 & 0 & 0 & 0 & \frac12 & 0 & 0\\
0 & 0 & 0 & 1 & \frac13 & \frac1{12} & \frac16 & 0\\
0 & 0 & -1 & 0 & -1 & -\frac32 & 0 & 1
\end{array}\right).
\]
因此，要求的一组辛基 $\{v_1,v_2,v_3,v_4\}$ 为
\[
(v_1,v_2,v_3,v_4)=(e_1,e_2,e_3,e_4)C,
\]
其中
\[
C=\left(\begin{array}{rrrr}
1 & 0 & \frac13 & -1\\
0 & \frac12 & \frac1{12} & -\frac32\\
0 & 0 & \frac16 & 0\\
0 & 0 & 0 & 1
\end{array}\right).
\]
$\square$

\noindent\textbf{例 10.22}\quad 设 $(V,g)$ 是辛空间，$\varphi$ 是 $V$ 上的辛变换，$\varphi$ 在一组辛基下的表示矩阵称为辛矩阵. 令 $A=\operatorname{diag}\{S,\cdots,S\}$，其中
\[
S=\left(\begin{array}{cc}
0 & 1\\
-1 & 0
\end{array}\right).
\]
求证：$n$ 阶方阵 $T$ 是辛矩阵的充要条件是 $T'AT=A$. 特别地，辛变换的行列式值等于 $1$ 或 $-1$.

\noindent\textbf{证明}\quad 选取一组辛基 $v_1,v_2,\cdots,v_n$，则交错型 $g$ 在这组基下的表示矩阵恰为 $A$. 设 $V$ 中向量 $x,y$ 在这组基下的坐标向量分别为 $\alpha,\beta$，辛变换 $\varphi$ 在这组基下的表示矩阵为 $T$，则
\[
\alpha'A\beta=g(x,y)=g(\varphi(x),\varphi(y))=(T\alpha)'A(T\beta)=\alpha'T'AT\beta
\]
对任意的列向量 $\alpha,\beta$ 都成立，从而 $T'AT=A$ 成立. 反之亦不难倒推回去，故结论成立. 由 $T'AT=A$ 取行列式可得 $|T|^2=1$，于是 $|T|=\pm1$，从而 $\det\varphi=\pm1$. $\square$

\noindent\textbf{注}\quad 由例 9.134 可知，实辛空间 $(V,g)$ 上的辛变换 $\varphi$ 的行列式值等于 $1$，请读者自行思考其中的原因.

\subsection*{§10.5\quad 对称型与正交几何}

\noindent\textbf{例 10.23}\quad 设 $g$ 是 $n$ 维线性空间 $V$ 上的非退化对称型，$\varphi$ 是 $V$ 上的正交变换，求证：$\det\varphi=\pm1$.

\noindent\textbf{证明}\quad 选取一组正交基 $e_1,e_2,\cdots,e_n$，则非退化对称型 $g$ 在这组基下的表示矩阵为可逆对角矩阵 $A$. 设 $\varphi$ 在这组基下的表示矩阵为 $T$，因为 $\varphi$ 是正交变换，故对任意的 $x,y\in V$ 有 $g(\varphi(x),\varphi(y))=g(x,y)$，由类似于上例的讨论可知 $T'AT=A$. 此式取行列式可得 $|T|^2=1$，于是 $|T|=\pm1$，从而 $\det\varphi=\pm1$. $\square$

\noindent\textbf{例 10.24}\quad 设 $g$ 是 $n$ 维线性空间 $V$ 上的非退化对称型，$\varphi$ 是 $V$ 上的线性变换，求证：$\varphi$ 是正交变换的充要条件是 $\varphi^*\varphi=I$.

\noindent\textbf{证明}\quad 若 $\varphi$ 是正交变换，则对任意的 $x,y\in V$ 有 $g(\varphi(x),\varphi(y))=g(x,y)$，从而可得
\[
g(x,\varphi^*\varphi(y))=g(x,y).
\]
因为 $g$ 非退化，所以 $\varphi^*\varphi=I$. 充分性只要反过来推回去即可. $\square$

\noindent\textbf{例 10.25}\quad 设 $V$ 是双曲平面，$\varphi$ 是 $V$ 上的正交变换且 $\det\varphi=-1$，求证：$\varphi$ 是镜像变换.

\noindent\textbf{证明}\quad 设 $g$ 是定义双曲平面 $V$ 的非退化对称型，则由教材 [1] 中的定理 10.5.3 可知，存在 $V$ 的一组基 $u,v$，使得
\[
g(u,u)=g(v,v)=0,\quad g(u,v)=g(v,u)=1.
\]
设 $\varphi$ 在基 $u,v$ 下的表示矩阵为 $A=(a_{ij})$，即有
\[
\varphi(u)=a_{11}u+a_{21}v,\quad \varphi(v)=a_{12}u+a_{22}v,
\]
则从 $g(\varphi(u),\varphi(u))=g(u,u)=0$ 可推出 $a_{11}=0$ 或 $a_{21}=0$；从 $g(\varphi(v),\varphi(v))=g(v,v)=0$ 可推出 $a_{22}=0$ 或 $a_{12}=0$. 由于 $\varphi$ 是可逆变换，故只有下列可能性：或者 $a_{12}=a_{21}=0$，或者 $a_{11}=a_{22}=0$. 假设为前者，由 $g(\varphi(u),\varphi(v))=g(u,v)=1$ 可推出 $a_{11}a_{22}=1$，这与 $\det\varphi=-1$ 矛盾，因此必有 $a_{11}=a_{22}=0$. 这时再由 $g(\varphi(u),\varphi(v))=g(u,v)=1$ 可推出 $a_{12}a_{21}=1$. 令 $\beta=a_{12}u-v$，则不难验证 $\varphi=S_\beta$ 是镜像变换. $\square$

\noindent\textbf{例 10.26}\quad 设 $g$ 是 $n$ 维实线性空间 $V$ 上的非退化对称型，$g$ 的正惯性指数为 $p$，负惯性指数为 $q$. 假设 $W$ 是 $V$ 的极大全迷向子空间，求证：$\dim W=\min\{p,q\}$.

\noindent\textbf{证明}\quad 设 $\{e_1,e_2,\cdots,e_n\}$ 是 $V$ 的一组基，使得 $g$ 在这组基下的表示矩阵为
\[
\operatorname{diag}\{1,\cdots,1,-1,\cdots,-1\},
\]
其中有 $p$ 个 $1$，$q$ 个 $-1$，不妨假设 $p\le q$（$p>q$ 类似可证）. 令
\[
v_1=e_1+e_{p+1},\quad v_2=e_2+e_{p+2},\quad \cdots,\quad v_p=e_p+e_{2p};
\]
\[
u_1=e_1-e_{p+1},\quad u_2=e_2-e_{p+2},\quad \cdots,\quad u_p=e_p-e_{2p}.
\]
显然，$\{v_1,\cdots,v_p,u_1,\cdots,u_p,e_{2p+1},\cdots,e_n\}$ 是 $V$ 的一组基. 由基向量 $\{v_1,\cdots,v_p\}$ 生成的子空间 $W$，其维数等于 $p$，且不难验证这是一个全迷向子空间.

现假设 $W$ 是维数大于 $p$ 的全迷向子空间，令 $U$ 是由基向量 $\{e_{p+1},\cdots,e_n\}$ 生成的子空间. 因为 $U$ 的维数为 $n-p$，$W$ 的维数大于 $p$，故由交和空间的维数公式可知 $U\cap W\ne0$. 任取 $0\ne\beta\in U\cap W$，可设 $\beta=b_{p+1}e_{p+1}+b_{p+2}e_{p+2}+\cdots+b_ne_n$，则
\[
g(\beta,\beta)=-b_{p+1}^2-b_{p+2}^2-\cdots-b_n^2<0,
\]
这与 $\beta$ 是迷向向量矛盾，因此 $V$ 中全迷向子空间的维数最多为 $p$. $\square$

\noindent\textbf{例 10.27}\quad 设 $V=U_1\oplus U_2\oplus\cdots\oplus U_r$，若 $U_i\perp U_j$ 对一切 $i\ne j$ 成立，则称 $V$ 是 $U_i$ 的正交直和，记为 $V=U_1\perp U_2\perp\cdots\perp U_r$. 现假设
\[
V=U_1\perp U_2\perp\cdots\perp U_r=V_1\perp V_2\perp\cdots\perp V_r,
\]
若存在保距同构 $\varphi_i:U_i\to V_i\ (1\le i\le r)$，求证：存在 $V$ 上的正交变换 $\varphi$，使之在每个 $U_i$ 上的限制就是 $\varphi_i$.

\noindent\textbf{证明}\quad 令 $\varphi(u_1+u_2+\cdots+u_r)=\varphi_1(u_1)+\varphi_2(u_2)+\cdots+\varphi_r(u_r)$，其中每个 $u_i\in U_i$. 不难验证 $\varphi$ 就是所需之正交变换. $\square$

\noindent\textbf{例 10.28}\quad 设 $g$ 是 $n$ 维实线性空间 $V$ 上的非退化对称型，$\varphi$ 是 $V$ 上的正交变换，$V_1=\{x\in V\mid \varphi(x)=x\}$. 求证：$V_1$ 是子空间且 $\dim V=\dim V_1+\dim(I-\varphi)V$.

\noindent\textbf{证明}\quad 显然 $V_1=\operatorname{Ker}(I-\varphi)$，由线性映射的维数公式即得结论. $\square$

\noindent\textbf{例 10.29}\quad 设 $V$ 是 $n$ 维欧氏空间，$g$ 是 $V$ 上的一个对称型. 满足 $g(u,u)=0$ 的非零向量 $u$ 称为迷向向量. 求证：$V$ 存在一组由迷向向量组成的标准正交基的充要条件是 $g$ 在 $V$ 的某一组标准正交基下的表示矩阵的迹等于零.

\noindent\textbf{证明}\quad 我们把问题转化成代数语言. 设 $A$ 是 $g$ 在某一组标准正交基下的表示矩阵，则 $A$ 是实对称矩阵. 因此原问题等价于下面的矩阵问题：实对称矩阵 $A=(a_{ij})$ 正交相似于主对角元全是零的对称矩阵的充要条件是 $\operatorname{tr}(A)=0$. 必要性是显然的，下证充分性. 对阶数 $n$ 进行归纳. 当 $n=1$ 时结论显然成立，假设对 $n-1$ 阶矩阵结论已成立，现证 $n$ 阶矩阵的情形. 下面分两种情况进行讨论.

若 $A$ 的主对角元中有一个为零，由于主对角元的对换是正交相似变换，故不妨设 $a_{11}=0$，于是
\[
A=\left(\begin{array}{cc}
0 & \alpha'\\
\alpha & A_1
\end{array}\right).
\]
注意到 $A_1$ 是 $n-1$ 阶实对称矩阵且 $\operatorname{tr}(A_1)=0$，故由归纳假设可知，存在 $n-1$ 阶正交矩阵 $R$，使得 $R'A_1R$ 的主对角元全为零. 令 $P=\operatorname{diag}\{1,R\}$，则 $P$ 是 $n$ 阶正交矩阵，使得 $P'AP$ 的主对角元全为零.

若 $A$ 的主对角元全部非零，我们的目标是通过正交相似变换将 $A$ 化为第 $(1,1)$ 元素为零的实对称矩阵，再由第一种情况的讨论即得结论. 首先设 $P$ 为正交矩阵，使得 $P'AP=\operatorname{diag}\{\lambda_1,\lambda_2,\cdots,\lambda_n\}$，由假设可知 $\lambda_1+\lambda_2+\cdots+\lambda_n=\operatorname{tr}(A)=0$. 若存在某个 $\lambda_i=0$，则结论得证，故不妨设 $\lambda_i$ 全部非零. 因为主对角元的对换是正交相似变换，故不妨进一步假设 $\lambda_1>0,\ \lambda_2<0$. 设 $P=(e_1,e_2,\cdots,e_n)$，则 $e_1,e_2,\cdots,e_n$ 是 $\mathbb{R}^n$ 的标准正交基，且满足 $Ae_i=\lambda_i e_i$. 设 $t,s$ 为实参数，满足如下条件：
\[
(e_1+te_2)'A(e_1+te_2)=\lambda_1+t^2\lambda_2=0,\quad
(e_1+te_2)'(e_1+se_2)=1+st=0,
\]
则容易算出
\[
t=\sqrt{-\frac{\lambda_1}{\lambda_2}},\quad
s=-\sqrt{-\frac{\lambda_2}{\lambda_1}}.
\]
令
\[
\widetilde e_1=\frac{e_1+te_2}{\|e_1+te_2\|},\quad
\widetilde e_2=\frac{e_1+se_2}{\|e_1+se_2\|},
\]
则 $\widetilde e_1,\widetilde e_2,e_3,\cdots,e_n$ 是 $\mathbb{R}^n$ 的标准正交基. 令 $Q=(\widetilde e_1,\widetilde e_2,e_3,\cdots,e_n)$，则 $Q$ 是正交矩阵，且由上述条件容易验证 $Q'AQ$ 是一个第 $(1,1)$ 元素为零的实对称矩阵，从而结论得证. $\square$

\noindent\textbf{例 10.30}\quad 设四维实空间 $V$ 上定义了一个对称型 $g$，在基 $\{e_1,e_2,e_3,e_4\}$ 下的表示矩阵为
\[
\left(\begin{array}{rrrr}
1 & 0 & 0 & 0\\
0 & 1 & 0 & 0\\
0 & 0 & 1 & 0\\
0 & 0 & 0 & -1
\end{array}\right).
\]
上述空间 $V$ 称为 Minkowski 空间. $V$ 中适合 $g(\alpha,\alpha)>0$ 的向量 $\alpha$ 称为空间向量；适合 $g(\alpha,\alpha)<0$ 的向量称为时间向量；适合 $g(\alpha,\alpha)=0$ 的非零向量称为光向量. 证明：

(1) $V$ 中任意两个时间向量不可能互相正交；

(2) $V$ 中任意一个时间向量不可能正交于一个光向量；

(3) $V$ 中两个光向量正交的充要条件是它们线性相关.

\noindent\textbf{证明}\quad (1) 设 $x,y$ 是两个时间向量且
\[
x=a_1e_1+a_2e_2+a_3e_3+a_4e_4,\quad
y=b_1e_1+b_2e_2+b_3e_3+b_4e_4.
\]
若它们正交，则 $g(x,y)=0$ 且 $g(x,x)<0,\ g(y,y)<0$. 于是
\[
a_1^2+a_2^2+a_3^2<a_4^2,\quad
b_1^2+b_2^2+b_3^2<b_4^2,\quad
a_1b_1+a_2b_2+a_3b_3=a_4b_4.
\]
由 Cauchy 不等式，有
\[
a_4^2b_4^2=(a_1b_1+a_2b_2+a_3b_3)^2
\le (a_1^2+a_2^2+a_3^2)(b_1^2+b_2^2+b_3^2),
\]
导出矛盾.

(2) 设 $u=c_1e_1+c_2e_2+c_3e_3+c_4e_4$ 是光向量，则 $g(u,u)=0$，即 $c_1^2+c_2^2+c_3^2=c_4^2$. 若 $x\perp u$，则 $a_1c_1+a_2c_2+a_3c_3=a_4c_4$. 同 (1) 用 Cauchy 不等式可证这是不可能的.

(3) 设 $v=d_1e_1+d_2e_2+d_3e_3+d_4e_4$ 也是光向量，若它和 $u$ 正交，则 $c_1d_1+c_2d_2+c_3d_3=c_4d_4$. 又 $d_1^2+d_2^2+d_3^2=d_4^2$，运用 Cauchy 不等式等号成立的充要条件即知 $u,v$ 线性相关. $\square$

\subsection*{§10.6\quad 基础训练}

\subsubsection*{10.6.1\quad 训练题}

\noindent\textbf{一、单选题}

1. 设 $\varphi$ 是有限维线性空间 $V\to U$ 的线性映射，$\varphi^*$ 是其对偶映射，则（\quad ）.

(A) 若 $\varphi$ 是单映射，则 $\varphi^*$ 也是单映射
\quad
(B) 若 $\varphi$ 是满映射，则 $\varphi^*$ 也是满映射

(C) 若 $\varphi$ 是单映射，则 $\varphi^*$ 是满映射
\quad
(D) 若 $\varphi^*$ 是单映射，则 $\varphi$ 也是单映射

2. 设 $V$ 是数域 $\mathbb{F}$ 上的三维空间，$V^*$ 是其对偶空间，$v_1,v_2,v_3$ 是 $V$ 的一组基，$v_1^*,v_2^*,v_3^*$ 是对偶基，则 $V$ 中基 $v_1,\ v_1+v_2,\ v_1+v_2+v_3$ 的对偶基是（\quad ）.

(A) $v_1^*,v_1^*+v_2^*,v_1^*+v_2^*+v_3^*$
\quad
(B) $v_1^*-v_2^*,v_2^*-v_3^*,v_3^*$

(C) $v_1^*,v_1^*-v_2^*,v_2^*-v_3^*$
\quad
(D) $v_1^*-v_2^*,v_1^*-v_3^*,v_2^*+v_3^*$

3. 设 $\varphi$ 是 $n$ 维线性空间 $V$ 上的线性变换，$\varphi^*$ 是其对偶变换. $V_1$ 是 $V$ 的子空间，记 $V_1^\perp=\{f\in V^*\mid f(V_1)=0\}$，则下列结论正确的是（\quad ）.

(A) $\dim(\operatorname{Ker}\varphi^*)^\perp=\dim\operatorname{Ker}\varphi$
\quad
(B) $\dim(\operatorname{Ker}\varphi)^\perp=\dim\operatorname{Im}\varphi$

(C) $\dim(\operatorname{Im}\varphi^*)^\perp=\dim\operatorname{Im}\varphi$
\quad
(D) $\dim\operatorname{Im}\varphi^*=\dim(\operatorname{Im}\varphi)^\perp$

4. 设 $V$ 是数域 $\mathbb{F}$ 上的三维行向量空间，$x=(x_1,x_2,x_3),\ y=(y_1,y_2,y_3)$，则下列函数是 $V\times V\to\mathbb{F}$ 的双线性函数的是（\quad ）.

(A) $f(x,y)=x_1^2+x_1y_1+2x_1y_2$

(B) $f(x,y)=x_1y_1+x_2y_3+2x_1-3y_3$

(C) $f(x,y)=x_1x_2+x_2y_3+2x_3y_3$

(D) $f(x,y)=2x_1y_2+x_2y_1-5x_3y_2+2x_3y_3$

5. 双线性型 $g:V\times V\to\mathbb{F}$ 在线性空间 $V$ 的不同基下的表示矩阵（\quad ）.

(A) 必相似\quad (B) 必合同\quad (C) 必正交相似\quad (D) 必相等

6. 设 $g:U\times V\to\mathbb{F}$ 是双线性型，$U^*$ 是 $U$ 的对偶空间，作映射 $\varphi:V\to U^*$，$\varphi(v)=g(-,v)$，则（\quad ）.

(A) $\operatorname{Ker}\varphi$ 为 $g$ 的右根子空间
\quad
(B) $\operatorname{Ker}\varphi$ 为 $g$ 的左根子空间

(C) $\operatorname{Im}\varphi$ 为 $g$ 的右根子空间
\quad
(D) $\operatorname{Im}\varphi$ 为 $g$ 的左根子空间

7. 下列纯量积非退化的是（\quad ）.

(A) $V$ 是 $\mathbb{F}$ 上的四维行向量空间，$x=(x_1,x_2,x_3,x_4),\ y=(y_1,y_2,y_3,y_4)$，$g(x,y)=x_1y_2+2x_1y_3-3x_3y_3$

(B) $V$ 是 $\mathbb{F}$ 上的四维行向量空间，$x=(x_1,x_2,x_3,x_4),\ y=(y_1,y_2,y_3,y_4)$，$g(x,y)=x_1y_1+2x_1y_3-2x_2y_1+x_2y_2+x_3y_2+4x_3y_3$

(C) $V$ 是 $n$ 维实列向量空间，$A$ 是 $n$ 阶幂等矩阵，$A\ne I_n$，$g(x,y)=x'Ay$

(D) $V$ 是 $n$ 阶矩阵组成的线性空间，$g(A,B)=\operatorname{tr}(AB)$

8. 设 $g,h$ 是 $V$ 上两个非退化的纯量积，$v_1,v_2,\cdots,v_n$ 是 $V$ 的基，$g$ 在这组基下的表示矩阵为 $A$，$h$ 在这组基下的表示矩阵为 $B$. 设 $\varphi$ 是 $V$ 上的线性变换，使得 $h(x,y)=g(\varphi(x),y)$，则 $\varphi$ 在上述基下的表示矩阵为（\quad ）.

(A) $A^{-1}B$\quad (B) $BA^{-1}$\quad (C) $B'(A^{-1})'$\quad (D) $(A^{-1})'B'$

9. 数域 $\mathbb{F}$ 上两个 $n$ 阶反对称矩阵合同的充要条件是（\quad ）.

(A) 它们相抵\quad (B) 它们的特征值相同\quad (C) 它们相似\quad (D) 作为复矩阵它们酉相似

10. $U,V$ 分别是数域 $\mathbb{F}$ 上的 $m$ 维和 $n$ 维线性空间，则由 $U\times V\to\mathbb{F}$ 全体双线性型组成的线性空间的维数是（\quad ）.

(A) $m$\quad (B) $n$\quad (C) $m+n$\quad (D) $mn$

\noindent\textbf{二、填空题}

1. 若 $\dim V=n$，则 $\dim V^*=$（\quad ）.

2. 设 $\varphi$ 是线性空间 $V$ 到 $U$ 的线性映射，它在 $V$ 和 $U$ 的一对基下的表示矩阵为 $A$，则 $\varphi$ 的对偶映射 $\varphi^*$ 在对偶基下的表示矩阵是（\quad ）.

3. 设 $V$ 和 $U$ 分别是数域 $\mathbb{F}$ 上的 $n$ 维和 $m$ 维线性空间，$\varphi$ 是 $V$ 到 $U$ 的线性映射. 已知 $\dim\operatorname{Im}\varphi=r$，则 $\dim\operatorname{Ker}\varphi^*=$（\quad ）.

4. 设 $V$ 是数域 $\mathbb{F}$ 上的 $n$ 维线性空间，$V^*$ 是 $V$ 的对偶空间，设 $\{e_1,\cdots,e_n\},\{e'_1,\cdots,e'_n\}$ 是 $V$ 的两组基，且从 $\{e_1,\cdots,e_n\}$ 到 $\{e'_1,\cdots,e'_n\}$ 的过渡矩阵是 $P$. 又假设 $\{f_1,\cdots,f_n\}$ 和 $\{f'_1,\cdots,f'_n\}$ 分别是 $\{e_1,\cdots,e_n\}$ 和 $\{e'_1,\cdots,e'_n\}$ 的对偶基，则从 $\{f_1,\cdots,f_n\}$ 到 $\{f'_1,\cdots,f'_n\}$ 的过渡矩阵是（\quad ）.

5. 设 $g$ 是 $V\times U\to\mathbb{F}$ 的双线性型，$g^*$ 是 $U^*\times V^*\to\mathbb{F}$ 的双线性型. 假设 $g$ 在 $V$ 的基 $v_1,v_2,\cdots,v_n$ 及 $U$ 的基 $u_1,u_2,\cdots,u_m$ 下的表示矩阵为 $A$，$g^*$ 在 $U^*$ 的对偶基 $u_1^*,u_2^*,\cdots,u_m^*$ 及 $V^*$ 的对偶基 $v_1^*,v_2^*,\cdots,v_n^*$ 下的表示矩阵为 $B$. 问 $A$ 和 $B$ 满足什么条件时，必有 $g(v_i,u_j)=g^*(u_j^*,v_i^*)$?

6. 设 $U$ 和 $V$ 分别是数域 $\mathbb{F}$ 上的 $m$ 维和 $n$ 维线性空间，$g$ 是 $U\times V\to\mathbb{F}$ 的双线性型. 假设 $g$ 在一对基下的表示矩阵为 $A$ 且 $A$ 的秩为 $r$，则 $g$ 的左根子空间的维数是（\quad ），$g$ 的右根子空间的维数为（\quad ）.

7. 设 $g$ 是线性空间 $V$ 上的纯量积，$\varphi$ 是 $V$ 上的线性变换，$v_1,v_2,\cdots,v_n$ 是 $V$ 的一组基，$g$ 在这组基下的表示矩阵为 $A$，$\varphi$ 在这组基下的表示矩阵为 $B$，则双线性型 $h(x,y)=g(\varphi(x),y)$ 在这组基下的表示矩阵为（\quad ）.

8. 设 $V$ 是由实数域上次数小于 $3$ 的多项式全体组成的线性空间，定义双线性型
\[
g(f_1(x),f_2(x))=\int_0^1 f_1(x)f_2(x)\,dx.
\]
写出 $g$ 在基 $1,x,x^2$ 下的表示矩阵.

9. 将 $V$ 上的纯量积 $g$ 表示为一个对称型和一个交错型之和.

10. 设 $V$ 是 $n$ 维线性空间，则 $V$ 上所有交错型组成的线性空间的维数为（\quad ）.

\subsubsection*{10.6.2\quad 训练题答案}

\noindent\textbf{一、单选题}

1. 应选择 (C). $\varphi$ 是单映射的充要条件是 $\varphi^*$ 是满映射.

2. 应选择 (B). 计算方法参考填空题 4.

3. 应选择 (B). 设 $\varphi$ 的秩为 $r$，则 $\dim\operatorname{Ker}\varphi=n-r$，故 $\dim(\operatorname{Ker}\varphi)^\perp=r=\dim\operatorname{Im}\varphi$.

4. 应选择 (D).

5. 应选择 (B).

6. 应选择 (A).

7. 应选择 (D). 事实上，(A)，(B) 中纯量积的表示矩阵都是奇异阵，因此是退化的. 当矩阵 $A$ 是幂等矩阵而非单位矩阵时，$A$ 必是奇异阵，故 (C) 也是退化的. 若 $g(AB)=\operatorname{tr}(AB)=0$ 对任意的矩阵 $B$ 成立，取 $B=E_{ji}$，则可得 $a_{ij}=0(1\le i,j\le n)$，即有 $A=0$，从而 $g$ 是非退化的.

8. 应选择 (D). 不妨设 $V$ 是 $n$ 维列向量空间，又设 $C$ 是 $\varphi$ 的表示矩阵，则
\[
h(x,y)=x'By=g(Cx,y)=(Cx)'Ay=x'C'Ay.
\]
因此 $C=(BA^{-1})'=(A^{-1})'B'$.

9. 应选择 (A). 同阶反对称矩阵合同的充要条件是它们的秩相等.

10. 应选择 (D). 选定 $U,V$ 的基后，$U\times V\to\mathbb{F}$ 的双线性型组成的线性空间和 $\mathbb{F}$ 上 $m\times n$ 矩阵组成的线性空间同构，因此维数等于 $mn$.

\noindent\textbf{二、填空题}

1. $n$.

2. $A'$.

3. 设 $A$ 是 $\varphi$ 的表示矩阵，由于 $A$ 和 $A'$ 的秩相等，故 $\operatorname{Im}\varphi$ 和 $\operatorname{Im}\varphi^*$ 具有相同的维数，因此 $\dim\operatorname{Ker}\varphi^*=m-r$.

4. 设 $P=(p_{ij})$，又设从 $\{f'_1,\cdots,f'_n\}$ 到 $\{f_1,\cdots,f_n\}$ 的过渡矩阵为 $Q=(q_{ij})$，于是
\[
f_j=q_{1j}f'_1+\cdots+q_{nj}f'_n,\quad
e'_i=p_{1i}e_1+\cdots+p_{ni}e_n,
\]
从而 $\langle f_j,e'_i\rangle=q_{ij}=p_{ji}$，即 $Q=P'$. 因此从 $\{f_1,\cdots,f_n\}$ 到 $\{f'_1,\cdots,f'_n\}$ 的过渡矩阵是 $(P')^{-1}$.

5. $B=A'$.

6. $g$ 的左根子空间的维数为 $m-r$，右根子空间的维数为 $n-r$.

7. $B'A$.

8. 表示矩阵为
\[
\left(\begin{array}{ccc}
1 & \frac12 & \frac13\\
\frac12 & \frac13 & \frac14\\
\frac13 & \frac14 & \frac15
\end{array}\right).
\]

9. $g=\dfrac12(g+g^*)+\dfrac12(g-g^*)$，其中 $g^*(x,y)=g(y,x)$.

10. 交错型的表示矩阵为反对称矩阵，因此答案为 $\dfrac12 n(n-1)$.

\end{document}
